\section{Precision of the fixed nearest-rounded positional bank}
\label{sec:precision-proofs}

The lower bound in this section concerns a particular observation, not
an arbitrary encoder with the same number of bits.  The earlier
finite-bank result already established a twelve-coordinate precision
floor together with noise and source-edit bounds \citep{BCJ001}.  Its floored
triple/half-count coordinates differ from the nearest-rounded direct
$J/C/D$ coordinates below.  Before quantization the two representations
are linearly interconvertible, but separate coordinatewise quantization
does not preserve that equivalence.  We therefore give a complete
lower proof for this changed cut.  Its common-cell and two-point
argument also includes $h=0$; it does not assert a new
statistical/quantization transition.

\subsection{The exact observation and result}

Fix integers $M\ge64$ and $h\ge0$, and set $Q=2^h$.
Let $\mathcal N_Q$ round to the nearest point of
$\{0,1/Q,\ldots,1\}$ using a fixed tie rule.
Use the counts in \eqref{eq:a-full-counts}--\eqref{eq:a-half-count},
and form the twelve-coordinate bank
\begin{equation}\label{eq:p-bank}
 \mathcal B_{M,h}(X)=
 \bigl(\mathcal N_Q(J_u/M),\mathcal N_Q(C_u/M),
       \mathcal N_Q(D_u/M)\bigr)_{u\in\{00,01,10,11\}}.
\end{equation}
The decoder observes exactly
\begin{equation}\label{eq:p-bank-cut}
 \mathcal O=(\mathcal B_{M,h}(X),E),\qquad
 E=(X_M,X_{M+1}),
\end{equation}
and source-independent public format metadata.  The endpoint is exact.
No raw history, scratch counts, other token banks, rounding residuals,
or additional feature port crosses this cut.

Let $\mathcal C_M$ be the clean class of homogeneous binary order-two
sources with arbitrary initial law admitting a stationary pair law
$\Pi$ with $P^2(u,v)\ge\Pi(v)/M$ for every $u,v$.  Define
\begin{equation}\label{eq:p-bank-minimax}
 R^*(M,h)=\inf_A\sup_{(\theta,\mu)\in\mathcal C_M}
 \E\,\TV\bigl(A(\mathcal O),\theta_{Y_M}\bigr).
\end{equation}
The infimum permits arbitrary computation, any binary probability
output, and arbitrary independent decoder seeds.  It is not restricted
to the gate decoder or to a finite output grid.  Public randomization
independent of the source and sample may be included in these seeds;
source-dependent advice or randomization correlated with hidden history
is not part of \eqref{eq:p-bank-cut}.

\begin{theorem}[Fixed-bank precision]
\label{thm:bank-precision}
For every integer $M\ge64$ and every integer $h\ge0$,
\begin{equation}\label{eq:p-bank-rate}
 \frac1{1600000}\max\{M^{-1/2},2^{-h}\}
 \le R^*(M,h)
 \le\min\bigl\{1,667160\max\{M^{-1/2},2^{-h}\}\bigr\}.
\end{equation}
Both lower components already hold on the stationary IID Bernoulli
subclass with parameter $p\in[1/8,1/4]$.  Consequently the same lower
bounds hold for an adversarial experiment that allows zero edits,
provided its decoder has no more information than the fixed cut
\eqref{eq:p-bank-cut}.
\end{theorem}

\subsection{Two elementary probability inequalities}

For completeness, we give the general concentration and testing
inequalities needed below.  They apply to finite input laws; randomized
decoder seeds will be handled separately without restricting their
support.

\begin{lemma}[Bounded differences]
\label{lem:p-bounded-differences}
Let $F$ be a real function of independent finite-valued coordinates,
and suppose changing coordinate $i$ alone changes $F$ by at most
$c_i\ge0$.  With $V=\sum_i c_i^2>0$, for every $t>0$,
\[
 \Prb(|F-\E F|\ge t)\le2\exp(-2t^2/V).
\]
If $V=0$, $F$ is constant on the input support.
\end{lemma}
\begin{proof}
Reveal the independent coordinates in order and let $Z_i$ be the
increments of the Doob martingale of $F$.  Conditional on any revealed
prefix, $Z_i$ has mean zero and lies in an interval of length at most
$c_i$.  Indeed, compare two values of coordinate $i$ by integrating
over the same independent remaining coordinates; the resulting
conditional expectations differ by at most $c_i$.

For a centered random variable $Z$ in an interval $[a,b]$ of length
$c$, put $g(\lambda)=\log\E e^{\lambda Z}$.  Its second derivative
is the variance under the exponentially tilted law.  For any law on
$[a,b]$, the pointwise inequality $(Z-a)(b-Z)\ge0$ implies
\[
 \Var(Z)\le(\E Z-a)(b-\E Z)\le(b-a)^2/4.
\]
The same bound therefore holds under every tilt.  Since
$g(0)=g'(0)=0$, integrating twice gives
$g(\lambda)\le\lambda^2c^2/8$ for either sign of $\lambda$.
Apply this bound conditionally to each $Z_i$ and iterate conditional
expectations to obtain
\[
 \E e^{\lambda(F-\E F)}\le
 \exp\!\left(\lambda^2\sum_i c_i^2/8\right).
\]
Markov's inequality, optimized at $\lambda=4t/V$, gives
$\Prb(F-\E F\ge t)\le e^{-2t^2/V}$.
Apply the same argument to $-F$ and add the two tails.
The zero-sensitivity case is immediate.
\end{proof}

\begin{lemma}[Pinsker's inequality for finite laws]
\label{lem:p-pinsker}
For finite probability laws $P,Q$,
\[
 \TV(P,Q)\le\sqrt{\KL(P\Vert Q)/2}.
\]
Moreover, relative entropy adds across independent product factors.
\end{lemma}
\begin{proof}
If $P$ has positive mass where $Q$ has zero mass, the divergence is
infinite and the bound is immediate.  Otherwise set
$A=\{x:P(x)>Q(x)\}$, so
$\TV(P,Q)=P(A)-Q(A)$.
For a subset with $P$-mass $a>0$ and $Q$-mass $b$, its contribution to
relative entropy is $a$ times an expectation of
$\log(P(x)/Q(x))$ under normalized $P$ weights.  Jensen's inequality
for the convex function $-\log$ bounds this below by
$a\log(a/b)$: indeed the normalized expectation of $Q(x)/P(x)$
is at most $b/a$, with equality if $Q$ places no mass on zeros of $P$
inside the subset.  Zero masses follow by continuity or infinite
divergence.  Apply this log-sum inequality to $A$ and its complement:
\[
 \KL(P\Vert Q)\ge\operatorname{kl}(P(A)\Vert Q(A)).
\]
For fixed $q\in(0,1)$, the binary divergence
$f(p)=\operatorname{kl}(p\Vert q)$ satisfies
$f(q)=f'(q)=0$ and $f''(p)=1/[p(1-p)]\ge4$.
Integrating twice from $q$ to $p$, in either direction, gives
$f(p)\ge2(p-q)^2$.  Boundary cases follow by continuity or infinite
divergence.  Thus $\KL(P\Vert Q)\ge2\TV(P,Q)^2$.
Finally, a product-density ratio is a product of the factor ratios;
expanding its logarithm into a sum and taking expectations proves
additivity for independent products.
\end{proof}

\subsection{A common rounding cell for the complete mean vector}

Consider stationary IID Bernoulli$(p)$ input symbols, with
$p\in I=[1/8,1/4]$, and put $t=\lfloor M/2\rfloor/M$.
The means of the unrounded normalized coordinates, in row order
$00,01,10,11$, are
\begin{align}
 \E_p(J_u/M)_u
   &=((1-p)^2,p(1-p),p(1-p),p^2),\label{eq:p-bank-J-means}\\
 \E_p(C_u/M)_u
   &=(p(1-p)^2,p^2(1-p),p(1-p)^2,p^2(1-p)),
     \label{eq:p-bank-C-means}\\
 \E_p(D_u/M)_u
   &=t((1-p)^2,p(1-p),p(1-p),p^2).
     \label{eq:p-bank-D-means}
\end{align}
Every summand has the displayed law, including the final triple ending
at $X_{M+1}$.  Overlapping pairs and triples need not be independent
for these identities, which use only linearity of expectation.
The half bank has $\lfloor M/2\rfloor$ summands and denominator $M$.
For $M\ge64$, $t\ge1/3$.

On $I$, the four $J$-coordinate derivative magnitudes are bounded below
by $3/2,1/2,1/2,1/4$, respectively.  The derivatives of the two
distinct $C$ polynomials are
\[
 (1-p)(1-3p)\ge3/16,\qquad
 p(2-3p)\ge13/64.
\]
The $D$ derivatives are $t$ times the corresponding $J$ derivatives.
Thus all twelve mean functions are strictly monotone on $I$, and every
derivative has magnitude at least $d=1/12$.  Duplicated mean functions
do not cause a problem, and none of the coordinates is constant.

The rounding thresholds are
\[
 \mathcal T_Q=\{(j+1/2)/Q:0\le j\le Q-1\}.
\]
Each monotone coordinate mean crosses at most $Q$ thresholds.
Take the distinct crossing points in the interior of $I$ and include
the two endpoints.  These partition $I$ into at most $12Q+1$ open
intervals.  One such interval $(a_*,b_*)$ has length
\begin{equation}\label{eq:p-cell-length}
 L=b_*-a_*\ge\frac{1/8}{12Q+1}\ge\frac1{104Q}.
\end{equation}
Choose $p_0=a_*+L/3$ and $p_1=a_*+2L/3$.  Then
\begin{equation}\label{eq:p-cell-separation}
 \delta_p:=p_1-p_0=L/3\ge\frac1{312Q},\qquad
 \delta_p\le\frac1{24}.
\end{equation}
The mean vectors at both parameters round to one common bank
$\mathbf b_*$.  More precisely, for each coordinate mean $f$ and
$j\in\{0,1\}$,
\begin{equation}\label{eq:p-cell-margin}
 \operatorname{dist}(f(p_j),\mathcal T_Q)
 \ge dL/3\ge\frac1{3744Q}=:\rho.
\end{equation}
To justify this at the first and last partition intervals as well,
no threshold lies strictly between the values $f(a_*)$ and $f(b_*)$.
A threshold outside this value interval is separated from $f(p_j)$
by at least its distance to the corresponding endpoint value; integrating
the derivative lower bound over the relevant subinterval gives at
least $dL/3$.  This also handles thresholds outside $f(I)$, without
extending the derivative bound beyond $I$.  Endpoint ties are harmless
because $p_0,p_1$ are interior.  The argument includes $h=0$, when
the sole threshold is $1/2$.

\subsection{Concentration of the actual bank and its joint endpoint law}

Let $F$ be any one of the twelve unrounded normalized coordinates.
Changing one raw bit affects at most three relevant departure triples,
so changes the count by at most three.  Pair and half-pair coordinates
in fact require at most two.  There are $N=M+2$ independent raw bits,
so the bounded-difference constants satisfy
\[
 \sum_i c_i^2\le(M+2)(3/M)^2\le10/M,\qquad M\ge64.
\]
Lemma~\ref{lem:p-bounded-differences} yields
\[
 \Prb_p(|F-\E_pF|\ge\rho)\le2\exp(-M\rho^2/5).
\]
If every coordinate is strictly within $\rho$ of its mean, the actual rounded
bank equals $\mathbf b_*$.  The union bound and
\eqref{eq:p-cell-margin} give, separately at $p_0$ and $p_1$,
\begin{equation}\label{eq:p-bank-concentration}
 \Prb_{p_j}(\mathcal B_{M,h}\ne\mathbf b_*)
 \le\epsilon:=24\exp\!\left(-\frac{M}{5\cdot3744^2Q^2}\right).
\end{equation}
This is concentration of the random bank, not a substitution of the
rounded expectation for the expectation of a rounded variable.
No independence between bank coordinates or count summands is needed.

Set $C_0=500000000$.  If $M\ge C_0Q^2$, then
$\epsilon\le1/16$.  For an elementary numerical verification,
$C_0>30\cdot3744^2$, and the nonnegative Taylor terms for $e^6$
through degree ten already exceed $384$.  Hence
$24e^{-C_0/(5\cdot3744^2)}<24/384=1/16$.

Let $P_j$ be the joint law of $(\mathcal B_{M,h},E)$ at $p_j$, and
let $\widetilde P_j$ be the law of $(\mathbf b_*,E)$ with the same
endpoint marginal.  Couple them by retaining exactly the same endpoint
$E$ and replacing only the bank.  Equation~\eqref{eq:p-bank-concentration}
implies
\[
 \TV(P_j,\widetilde P_j)\le\epsilon.
\]
Under IID sampling, the endpoint marginal is
$\operatorname{Bern}(p_j)^{\otimes2}$.  Couple the two endpoint bits
using common independent uniforms at the two parameters.  A union
bound gives $\TV(\widetilde P_0,\widetilde P_1)\le2\delta_p$.
Therefore, when $M\ge C_0Q^2$,
\begin{equation}\label{eq:p-bank-TV}
 \TV(P_0,P_1)\le2\epsilon+2\delta_p
 \le\frac18+\frac1{12}=\frac5{24}<\frac12.
\end{equation}
The bank is not asserted independent of its endpoint: the final
departures explicitly use the protected bits.  Keeping the same
endpoint in the first coupling is what validates this argument.

\subsection{Randomized two-point loss and the full-history floor}

For any decoder let $A$ denote its probability assigned to symbol one.
In the IID submodel the true probability is $p$ at every actual
endpoint, so total variation loss is exactly $|A-p|$.
For two finite observation laws $P_0,P_1$, their common
subprobability measure has mass $1-\TV(P_0,P_1)$.
Append the same independent seed law to both observations, allowing
arbitrary seed spaces.  This preserves total variation and gives a
common subprobability product measure of the same mass.  On every
observation and seed,
$|A-p_0|+|A-p_1|\ge|p_1-p_0|$.
Integrating this inequality against that common measure, and using
nonnegativity to bound the two risks below, proves
\begin{equation}\label{eq:p-two-point}
 \max_{j=0,1}\E_j|A-p_j|
 \ge\frac{|p_1-p_0|}{2}\bigl(1-\TV(P_0,P_1)\bigr).
\end{equation}
This covers arbitrary independent randomized decoders and public
source-independent randomization.  Apply \eqref{eq:p-bank-TV} and
\eqref{eq:p-cell-separation} to obtain
\begin{equation}\label{eq:p-large-regime}
 R^*(M,h)\ge\frac{\delta_p}{4}\ge\frac1{1248Q},
 \qquad M\ge C_0Q^2.
\end{equation}

For the complementary regime, we prove a floor even when the decoder
sees all $N=M+2$ raw bits.  Take
$p_0=1/8$ and $p_1=p_0+\delta$ with
$\delta=1/(16\sqrt N)$, so both belong to $I$.
The inequality $\log x\le x-1$ gives the Bernoulli divergence bound
\[
 \begin{split}
 \KL(\operatorname{Bern}(p_0)\Vert\operatorname{Bern}(p_1))
 &\le p_0\left(\frac{p_0}{p_1}-1\right)
 +(1-p_0)\left(\frac{1-p_0}{1-p_1}-1\right)\\
 &=\frac{\delta^2}{p_1(1-p_1)}
 \le\frac{64}{7}\delta^2.
 \end{split}
\]
The last bound uses $p_1\in I$, where $p_1(1-p_1)\ge7/64$.
Product additivity in Lemma~\ref{lem:p-pinsker} bounds the full-history
divergence by $N(64/7)\delta^2=1/28$, and therefore its total variation
by $\sqrt{1/56}<1/4$.  Passing to any function of the history cannot
increase total variation: each event in the output space has a
preimage event in the history space.  Independent randomization also
cannot increase it, by integration over the common seed law.
Consequently \eqref{eq:p-two-point} proves
\begin{equation}\label{eq:p-history-floor}
 R^*(M,h)\ge\frac{3\delta}{8}
 =\frac3{128\sqrt{M+2}}\ge\frac1{64\sqrt M}.
\end{equation}
The last inequality is equivalent to $9M\ge4(M+2)$, valid here.
These IID tables are legitimate all-state witnesses: put
$\theta_u(1)=p$ for every row and initialize the pair with
$\Pi=\operatorname{Bern}(p)^{\otimes2}$.  After two transitions,
$P^2(u,v)=\Pi(v)$ for every initial pair $u$, so the required
certificate $1/M$ holds on every row.

\subsection{Combining the regimes and interpreting precision}

\begin{proof}[Proof of Theorem~\ref{thm:bank-precision}]
For $M<C_0Q^2$, \eqref{eq:p-history-floor} gives
\[
 R^*(M,h)>\frac1{64\sqrt{C_0}\,Q}
 \ge\frac1{1600000Q},
\]
since $\sqrt{C_0}<25000$.
For $M\ge C_0Q^2$, \eqref{eq:p-large-regime} gives this same weaker
lower bound.  Together with \eqref{eq:p-history-floor}, both of
which hold for every decoder, this proves
\[
 R^*(M,h)\ge\frac1{1600000}
              \max\{M^{-1/2},Q^{-1}\}.
\]
The parameter pair used for the statistical floor need not equal the
pair used for the precision floor.  Each is a lower bound for every
decoder's source supremum, so their maximum is a lower bound for the
same minimax risk.

For the upper bound, apply \eqref{eq:a-rounded-risk} with $k=0$ and
$\xi=0$ to exactly the bank in \eqref{eq:p-bank}.  Its sequential
clipping and ratio gate are permitted arbitrary decoders, and
\[
 R^*(M,h)\le\min\left\{1,
 \frac{131072}{\sqrt M}+\frac{536088}{Q}\right\}
 \le\min\left\{1,667160\max\{M^{-1/2},Q^{-1}\}\right\}.
\]
This proves the complete finite-domain statement.
\end{proof}

For any fixed $K>0$, attaining $R^*(M,h)\le K/\sqrt M$ requires
\[
 Q\ge\frac{\sqrt M}{1600000K},\qquad
 h\ge\tfrac12\log_2 M-\log_2(1600000K).
\]
Conversely, $h\ge\lceil\tfrac12\log_2 M\rceil$ suffices by
\eqref{eq:a-rounded-risk}; a separate finite output stage with error
at most $2^{-b}$ adds just that error.  The necessity statement is
about fractional precision of this deterministic nearest-rounded
$J/C/D$ bank.  It is not a lower bound for arbitrary encoders using
$12(h+1)$ bits, arbitrary positional summaries, neural architectures,
or total arithmetic precision.  It gives neither a noise lower bound
nor a $k/M$ lower bound for this changed bank.  Accumulation precision,
addresses, raw input storage, score precision, parameters, masks, and
other program resources are not lower bounded here.  Additional hidden
history or counts would change the premise.  The deliberately loose
crossover $C_0$ is a proof device, not a practical sample-size threshold.
